\chapter{Conclusion générale}
\label{ch:conclusion}

Le PFA a consisté à concevoir, réaliser et \textbf{mettre en production} LexHub (\url{https://lexhub.space}) chez WebVoice, sous tutelle de l'ENSA Tétouan : un SI de cabinet dont le cœur est le couple dossier + plaidoirie, et non un chatbot autonome.

L'atelier en cinq étapes, la clarification QCM, l'assistant de révision (réponse + édition), la règle « assistant = plaidoiries validées », l'isolation multi-tenant et le modèle d'abonnement sont en service. L'architecture (navigateur $\rightarrow$ Caddy $\rightarrow$ FastAPI $\rightarrow$ PostgreSQL / Redis / MinIO, plus services cloud STT, OCR et LLM) et l'urbanisation MADI (chapitre~\ref{ch:conception}) cadrent un SaaS exploitable.

Le paiement carte opérateur reste à activer ; les perspectives (section~\ref{sec:perspectives}) prolongent ce positionnement sans s'y substituer. LexHub accélère le passage de la parole à l'acte : le cabinet reste maître de ses dossiers et responsable de ce qu'il plaide.
